\documentclass[11pt,a4paper]{article}
\usepackage[left=22mm,right=22mm,top=21mm,bottom=22mm,headheight=15pt]{geometry}
\usepackage[T1]{fontenc}
\usepackage[spanish,es-noshorthands,es-nodecimaldot]{babel}
\usepackage{tgpagella}
\usepackage[scale=.94]{tgheros}
\usepackage{inconsolata}
\usepackage{microtype}
\usepackage{graphicx,booktabs,tabularx,array,multicol}
\usepackage{amsmath,amssymb}
\usepackage{xcolor}
\usepackage{enumitem}
\usepackage{titlesec,fancyhdr,lastpage}
\usepackage{caption}
\usepackage{tikz}
\usetikzlibrary{arrows.meta,positioning}
\usepackage{hyperref}
\usepackage{bookmark}

\definecolor{ink}{HTML}{172D46}
\definecolor{teal}{HTML}{167D8D}
\definecolor{amber}{HTML}{B67518}
\definecolor{muted}{HTML}{536579}
\definecolor{light}{HTML}{EEF3F6}
\hypersetup{colorlinks=true,linkcolor=teal,urlcolor=teal,citecolor=teal,
  pdftitle={Actividad de escritura, acceso factual y diseño de controles en HyTorch},
  pdfauthor={Hyphae Research Foundation},
  pdfsubject={Informe técnico V5/V6. Resultados medidos y protocolo prospectivo},
  pdfkeywords={HyTorch, acceso factual, control experimental, reproducibilidad, V5, V6}}
\urlstyle{tt}
\setlength{\parindent}{0pt}
\setlength{\parskip}{5pt}
\setlength{\emergencystretch}{1em}
\setlength{\tabcolsep}{5pt}
\renewcommand{\arraystretch}{1.16}
\setlist[itemize]{leftmargin=4.5mm,itemsep=3pt,topsep=3pt}
\setlist[enumerate]{leftmargin=6mm,itemsep=3pt,topsep=3pt}
\titleformat{\section}{\sffamily\Large\bfseries\color{ink}}{\thesection}{.65em}{}
\titleformat{\subsection}{\sffamily\normalsize\bfseries\color{ink}}{\thesubsection}{.55em}{}
\titlespacing*{\section}{0pt}{0pt}{10pt}
\titlespacing*{\subsection}{0pt}{10pt}{3pt}
\captionsetup{font=small,labelfont={bf,sf,color=ink},textfont=normalfont,skip=7pt}
\pagestyle{fancy}
\fancyhf{}
\fancyhead[L]{\sffamily\footnotesize\color{muted}HYPHAE RESEARCH FOUNDATION / HYTORCH}
\fancyhead[R]{\sffamily\footnotesize\color{muted}INFORME TÉCNICO 01}
\fancyfoot[L]{\sffamily\footnotesize\color{muted}12 de septiembre de 2026 \enspace / \enspace Versión 1.0}
\fancyfoot[R]{\sffamily\footnotesize\color{muted}\thepage\ / \pageref{LastPage}}
\renewcommand{\headrulewidth}{.3pt}
\renewcommand{\footrulewidth}{0pt}
\newcommand{\sys}{\textsc{HyTorch}}
\newcommand{\status}{\texttt{no-control-match}}
\newcommand{\factloss}{\Delta_{\mathrm{Value}}}
\newcommand{\callout}[2]{\par\vspace{3pt}\begingroup\setlength{\fboxsep}{9pt}%
  \colorbox{light}{\begin{minipage}{\dimexpr\linewidth-2\fboxsep\relax}%
  {\sffamily\bfseries\color{ink}#1}\par\small #2\end{minipage}}\endgroup\par\vspace{3pt}}
\newcommand{\figurepanel}[3]{\begin{center}\includegraphics[width=\linewidth]{figures/#1.pdf}%
  \captionof{figure}{#2}\label{#3}\end{center}}
\newcommand{\hash}[1]{{\small\ttfamily #1}}
\newcolumntype{Y}{>{\raggedright\arraybackslash}X}
\newcommand{\pageend}{\clearpage}
\begin{document}
\thispagestyle{empty}

{\sffamily\bfseries\color{teal}\small HYPHAE RESEARCH FOUNDATION\hfill INFORME TÉCNICO 01}
\vspace{15mm}

{\sffamily\bfseries\color{ink}\fontsize{29}{34}\selectfont
Actividad de escritura,\\
acceso factual y\\
diseño de controles\\
en HyTorch\par}
\vspace{5mm}
{\sffamily\large\color{muted}Resultados V5 y protocolo prospectivo V6}

\vspace{6mm}
{\color{teal}\rule{\linewidth}{1.2pt}}
\vspace{2mm}

\begin{tabularx}{\linewidth}{@{}YYY@{}}
{\sffamily\bfseries\Large 2.190} & {\sffamily\bfseries\Large 50 min 46 s} & {\sffamily\bfseries\Large 2 baterías}\\[-2pt]
{\small lotes de medición completos} & {\small ejecución y cierre supervisado} & {\small coincidencia numérica exacta}\\
\end{tabularx}

\vspace{6mm}
\textbf{Resumen.} Un modelo de \sys{} entrenado en un mundo factual sintético
alcanza el criterio de competencia en hechos frecuentes de descubrimiento. Al vetar por separado las
escrituras de sus cuatro unidades, la pérdida del valor factual y la pérdida
del texto original responden de forma distinta. U0 produce el mayor deterioro
factual medio; U1 tiene un efecto factual medio pequeño y heterogéneo, pese a
alterar la pérdida del texto de control.

La medición completa superó sus comprobaciones de integridad y repetibilidad.
Sin embargo, ninguna de las otras unidades satisface todos los criterios fijados
para servir como control comparable de U0. El resultado final es \status.

\callout{Conclusión y alcance}{
Hay evidencia de contribución al \emph{acceso factual} en el modelo ya entrenado.
Todavía no hay una demostración de que un canal deje de aprender durante el
entrenamiento y provoque alucinaciones. El protocolo V6 define una comparación
nueva entre bloqueo persistente y distribuido; sus ramas no se han entrenado.}

\vfill
{\small\color{muted}
Documento de resultados y diseño experimental. No es un artículo de investigación
ni una afirmación confirmatoria. Incluye datos reducidos, gráficas vectoriales,
código de reconstrucción y referencias de procedencia.\\[4pt]
12 de septiembre de 2026 \quad / \quad Versión 1.0\\
\url{https://github.com/Hyphae-Research-Foundation/hytorch}\\[4pt]
Documento, figuras y datos: CC BY-SA 4.0. Código: Apache-2.0.\\
Verificaciones de software y revisión interna; V6 no tiene validación empírica.}

\pageend
\section{Pregunta, sistema y alcance}
\label{sec:scope}

La hipótesis que motiva el trabajo es que una función factual podría dejar de
mejorar mientras otras rutas sostienen una pérdida agregada favorable. Para
estudiarla deben separarse tres afirmaciones:

\begin{enumerate}
\item Una métrica agregada puede ocultar diferencias de desempeño entre funciones.
\item Un fallo local durante el aprendizaje contribuye causalmente a errores
factuales y, bajo ciertas condiciones de respuesta, a afirmaciones falsas.
\item El mecanismo es frecuente en otros modelos y explica una parte relevante
de sus alucinaciones.
\end{enumerate}

La evidencia de este informe informa sobre la primera pregunta y prepara una
intervención para la segunda. No estima la frecuencia del mecanismo en modelos
ordinarios ni su contribución a todas las alucinaciones.

\subsection{Qué es una unidad y qué se bloquea}
El modelo contiene cuatro sitios de escritura, U0--U3, uno por bloque. Las
propuestas pasan por validación, asignación y aplicación cuantizada al residual.
El veto experimental convierte los COMMIT sobrevivientes de una unidad en
denegaciones explícitas, manteniendo los frames, los abortos originales y la
cuantización de la ruta de cero commits. No se promueven candidatos perdedores.

Una escritura activa, un gradiente no nulo o un codebook que cambia no acreditan
conocimiento factual por sí solos. A la inversa, un efecto pequeño bajo una
ablación puede reflejar redundancia o variación entre hechos; no identifica
automáticamente una unidad que no aprendió.

\subsection{Tarea y referencia}
El mundo sintético asigna valores a entidades y relaciones. Se distinguen hechos
frecuentes, raros y no expuestos, además de distintas frecuencias de mención de
entidades. La verdad de una respuesta procede del oráculo del mundo: una respuesta
correcta por azar a un hecho no expuesto continúa siendo correcta.

La referencia V5 sigue un warmup auxiliar W640 y 4.000 updates de adquisición.
Usa secuencias de longitud 128, vocabulario de 428 tokens y batches de 48 filas:
32 originales, ocho auxiliares y ocho canónicas. El objetivo combina las tres
pérdidas con pesos $1$, $0{,}25$, $0{,}25$ y una regularización global de peso
$0{,}001$. Esta receta describe la referencia evaluada; no permite inferir su
calidad a partir del valor de la loss solamente.

\begin{tabularx}{\linewidth}{@{}p{32mm}Y@{}}
\toprule
\textbf{Estado} & \textbf{Qué existe al cierre de este informe}\\
\midrule
Observado & Competencia de la referencia y medición de los cuatro sitios en dos baterías completas.\\
Verificado & Capturas, reducciones, registros, cierre y selección bajo la regla fijada.\\
Diseñado & Panel V6 de 17 ramas y calendarios completos con pruebas mecánicas.\\
Pendiente & Calificación de la ejecución V6, entrenamiento de sus ramas y resultados causales.\\
\bottomrule
\end{tabularx}

\pageend
\section{Competencia factual de la referencia}
\label{sec:competence}

La referencia satisface el gate de competencia \emph{common-discovery}: al menos
80\% de exactitud macro y superioridad estricta frente a los cuatro priors
declarados en cada mapa. Este resultado corresponde al subconjunto de hechos
frecuentes. El desempeño sobre la población completa es considerablemente menor.

\figurepanel{competence}{Exactitud de la referencia por estrato y mapa de posiciones.
Los denominadores se conservan en cada serie. La exactitud de hechos frecuentes
no representa la exactitud de las 128 preguntas. No se presentan intervalos de
confianza: mapas y preguntas no son réplicas independientes de entrenamiento.}{fig:competence}

\begin{tabularx}{\linewidth}{@{}Yrrrr@{}}
\toprule
\textbf{Mapa} & \textbf{Frecuentes} & \textbf{Raros} & \textbf{No expuestos} & \textbf{Total}\\
\midrule
Identidad & 38/38 & 26/38 & 3/52 & 67/128\\
Desplazamiento 8 & 36/38 & 30/38 & 4/52 & 70/128\\
Separaciones 4 & 37/38 & 23/38 & 3/52 & 63/128\\
\bottomrule
\end{tabularx}

Se evaluaron 64 hechos de descubrimiento con dos preguntas por hecho: 19
frecuentes, 19 raros y 26 no expuestos. En el conjunto completo, las exactitudes
son 52,34\%, 54,69\% y 49,22\%. En frecuentes son 100\%, 94,74\% y 97,37\%.
Los cuatro priors comunes alcanzan 0\%, 5,263\%, 10,526\% y 10,526\%.

Las 384 respuestas de esta evaluación están presentes. La cobertura es 100\%,
sin abstenciones ni salidas inválidas; los errores son afirmaciones falsas en
este benchmark. Estos resultados pertenecen a la referencia ya entrenada y
no comparan ramas causales de bloqueo o rescate.

\callout{Una competencia suficiente, de alcance acotado}{
La referencia permite estudiar contribuciones al acceso en hechos frecuentes
que el sistema puede resolver. No acredita conocimiento uniforme de hechos raros
o no expuestos, ni capacidad lingüística general.}

\pageend
\section{Cómo se realizó la medición}
\label{sec:method}

Se congeló la referencia A4000 y se ejecutaron dos baterías con pertenencia,
orden y reglas previamente fijados. Cada input se evaluó sin veto y con cada
uno de los cuatro vetos singleton. No hubo actualizaciones de parámetros ni
entrenamiento durante esta fase.

\begin{center}
\begin{tikzpicture}[node distance=5mm,box/.style={draw=teal!45,fill=light,rounded corners=2pt,
  text width=33mm,minimum height=16mm,align=center,font=\sffamily\small},
  arr/.style={-{Stealth[length=2mm]},draw=teal,thick}]
\node[box] (a) {Referencia e inputs\\fijados};
\node[box,right=of a] (b) {Sano + cuatro vetos\\dos baterías};
\node[box,right=of b] (c) {Captura y lectura\\independiente};
\node[box,right=of c] (d) {Regla de selección\\sin modificar};
\draw[arr] (a)--(b);\draw[arr] (b)--(c);\draw[arr] (c)--(d);
\end{tikzpicture}
\end{center}

\subsection{Localización factual}
La pérdida factual se mide en la frontera pública de respuesta, antes de
suministrar el token Value correcto. Se usa la distribución completa de 428
clases, sin renormalizar sobre el conjunto de posibles valores:
\begin{align}
\ell_{f,q,p}^{(u)}&=-\log\operatorname{softmax}(z_{f,q,p}^{(u)})[v_f],\\
\factloss(u)&=\frac{1}{|F|}\sum_{f\in F}
  \frac{1}{|Q_f||P|}\sum_{q\in Q_f,p\in P}
  \left(\ell_{f,q,p}^{(u)}-\ell_{f,q,p}^{(\varnothing)}\right).
\end{align}
Cada hecho tiene el mismo peso; primero se promedian sus preguntas y posiciones.
Un valor positivo significa que vetar la unidad reduce, en media geométrica,
la probabilidad asignada al valor correcto. Los nats no son puntos porcentuales de exactitud.

\subsection{Control de perturbación general}
El inventario contiene las 1.120 ocurrencias originales de texto de entidades
discovery, manteniendo multiplicidades, mezcla factual/mención, targets válidos
y EOS. La NLL se pondera por tokens dentro de cada mapa y los tres mapas reciben
igual peso. Para cada unidad se conservan:
\begin{tabularx}{\linewidth}{@{}p{15mm}Y@{}}
\toprule
$g_u$ & Cambio firmado de NLL de texto original: veto menos sano.\\
$m_u$ & RMS de la contribución local aplicada, respecto de la salida cuantizada de cero commits.\\
$r_u$ & RMS de esa salida local de cero commits; $\rho_u=m_u/r_u$.\\
$p_u$ & Candidatos admitidos entre todos los candidatos sanos válidos, incluidos overflow y abortos en el denominador.\\
\bottomrule
\end{tabularx}

Las RMS se obtienen de sumas de cuadrados y conteos globales, no de promediar RMS
de batches. El proxy $g$ usa texto conocido de esta tarea; no mide equivalencia
funcional de dos unidades ni perturbación de todo el lenguaje.

La población completa comprende 219 conjuntos de inputs, 2.190 lotes, 34.740
coordenadas reales y 300 filas dummy excluidas de las métricas. No se filtraron
hechos según los aciertos de la referencia.

\pageend
\section{Contribución de las cuatro unidades}
\label{sec:units}

U0 encabeza el localizador en ambas baterías: $\factloss=0{,}880215$ nats, con
una diferencia de $0{,}241685$ respecto de U2. U2 y U3 también muestran efectos
factuales medios positivos. La contribución observada no es exclusiva de U0.

\figurepanel{channel-effects}{Efectos factuales y de texto original, presentados con
sus escalas y desagregación por posición. Las métricas responden a poblaciones y
denominadores distintos. Ambas baterías producen los mismos valores; dibujar dos
copias no añadiría una réplica independiente.}{fig:units}

\begin{tabularx}{\linewidth}{@{}Yrrrrrr@{}}
\toprule
\textbf{Unidad} & $\factloss$ & $g$ & $m$ & $r$ & $\rho$ & $p$\\
\midrule
U0 & 0,880215 & 0,114421 & 0,533184 & 1,178609 & 0,452384 & 0,641188\\
U1 & $-0{,}021274$ & 0,142184 & 0,528574 & 1,293221 & 0,408727 & 0,637110\\
U2 & 0,638529 & 0,042236 & 0,541165 & 1,417010 & 0,381906 & 0,674473\\
U3 & 0,504185 & 0,030009 & 0,540632 & 1,529304 & 0,353515 & 0,658546\\
\bottomrule
\end{tabularx}
{\footnotesize Valores redondeados para lectura. Las dos baterías completas y los
valores de precisión completa se incluyen en los datos del informe.}

U1 presenta un efecto factual medio pequeño, ligeramente negativo, junto a un
incremento de NLL en texto original mayor que el de U0. Esta diferencia agregada
es informativa, pero depende de los hechos y mapas evaluados. Sus escrituras
están activas: el resultado no permite llamarlo un canal muerto o un módulo
dedicado exclusivamente al lenguaje general.

\pageend
\section{Por qué no se obtuvo un control emparejado}
\label{sec:matching}

La regla exigía un efecto factual mayor que $0{,}05$ nats en ambas baterías,
selección estable y un control distinto del objetivo que satisficiera cuatro
criterios. Para el líder U0, el margen de pérdida de texto es:
\[
s_g=\max(0{,}02;\ 0{,}20|g_{U0}|)=0{,}0228841794.
\]
Por tanto, un candidato debía tener $g$ dentro de
$[0{,}091536718;\ 0{,}137305076]$.

\figurepanel{matching}{Evaluación de los candidatos con los límites fijados. La banda
representa un criterio de admisibilidad, no un intervalo de confianza. Los tres
candidatos incumplen la condición de NLL de texto; pasan las otras tres.}{fig:matching}

\begin{tabularx}{\linewidth}{@{}Yrrrr@{}}
\toprule
\textbf{Candidato} & $|g_u-g_{U0}|$ & $m_u/m_{U0}$ & $\rho_u/\rho_{U0}$ & $|p_u-p_{U0}|$\\
\midrule
U1 & 0,027763 & 0,991354 & 0,903495 & 0,004078\\
U2 & 0,072185 & 1,014968 & 0,844207 & 0,033285\\
U3 & 0,084411 & 1,013968 & 0,781448 & 0,017358\\
\midrule
Límite & $\leq0{,}022884$ & $[2/3;1{,}5]$ & $[2/3;1{,}5]$ & $\leq0{,}10$\\
\bottomrule
\end{tabularx}

U1 es el candidato más cercano y aun así supera el límite de NLL en
$0{,}004878687$ nats. El resultado no depende del redondeo: la decisión usa los
valores completos. No se ampliaron calipers, no se sustituyó el objetivo y no
se añadieron baterías para obtener un resultado favorable.

\textbf{Resultado final: \status.} Los identificadores de objetivo y control
quedan nulos. U0 conserva su condición de líder del localizador; la comparación
emparejada original no queda admitida.

\pageend
\section{Heterogeneidad: el promedio no describe cada hecho}
\label{sec:heterogeneity}

La desagregación por hecho conserva información que se pierde en una sola barra.
Los 19 efectos medios por hecho de U0 son positivos, pero sus magnitudes difieren.
En U1, ocho hechos tienen efecto positivo y once negativo. Su rango va
aproximadamente de $-0{,}249738$ a $0{,}237074$ nats.

\figurepanel{fact-heterogeneity}{Efecto factual por hecho y unidad, promediado sobre
dos preguntas y tres mapas. Se conservan los 19 hechos y una escala centrada en
cero. Valores redondeados: los muy pequeños pueden aparecer como cero. Las
celdas no son réplicas independientes de entrenamiento.}{fig:heterogeneity}

\begin{tabularx}{\linewidth}{@{}Yrrrr@{}}
\toprule
\textbf{Mapa} & $\factloss(U0)$ & $\factloss(U1)$ & $g_{U0}$ & $g_{U1}$\\
\midrule
Identidad & 1,194326 & $-0{,}050548$ & 0,148076 & 0,050147\\
Desplazamiento 8 & 0,652737 & $-0{,}024324$ & 0,119203 & 0,172679\\
Separaciones 4 & 0,793581 & 0,011050 & 0,075983 & 0,203725\\
\bottomrule
\end{tabularx}

La mayor perturbación de texto de U1 frente a U0 en el promedio no ocurre en
identidad: aparece con las transformaciones de posición. Tampoco el signo del
efecto factual de U1 es constante entre los tres mapas. Por ello, el resultado
agregado no identifica una separación simple entre una unidad de ``hechos'' y
otra de ``lenguaje''.

Esta heterogeneidad no explica por sí sola un fallo de aprendizaje persistente.
Esa hipótesis requiere comparar trayectorias de entrenamiento intervenidas
con sus controles.

\pageend
\section{Integridad, repetibilidad y límites de interpretación}
\label{sec:integrity}

La medición cerró en 3.046,425 segundos, dentro de los 90 minutos autorizados,
sin reinicio automático. El supervisor registró retorno cero, grupo de procesos
vacío y cierre confirmado. El resultado final dejó de ser provisional antes
de producir las tablas de este informe.

\begin{tabularx}{\linewidth}{@{}p{47mm}Y@{}}
\toprule
\textbf{Comprobación} & \textbf{Resultado y alcance}\\
\midrule
Captura completa & 2.190 lotes y ambas baterías; ninguna coordenada omitida para mejorar la selección.\\
Lectura independiente & 19.656.415.525 bytes leídos; reducciones numéricas y tablas reconstruidas desde las capturas.\\
Registros Native & 2.193 registros ordenados comprobados: inicio de run y compromisos de inicio, lotes y fin. Consulta sobre copia nueva del ledger detenido.\\
Custodia final & Continuidad de namespace y metadatos de 39.722 archivos en 4.383 directorios.\\
Repetición & Deriva máxima NLL y RMS: cero; conteos enteros de clasificación idénticos.\\
Selección & Regla y pertenencia completas recalculadas; \status\ en el cierre.\\
\bottomrule
\end{tabularx}

\subsection{Qué acreditan estas comprobaciones}
La lectura raw verifica sus reducciones; Native enlaza registros y recibos,
pero no calcula por sí mismo la verdad factual ni autentica toda la aritmética
neuronal. La comprobación final de custodia conserva la continuidad de archivos
ya hasheados en un filesystem confiable; no es una garantía frente a reescrituras
adversarias que preserven metadatos. Los controles tienen alcances complementarios.

La ruta B16 se calificó contra la referencia y el sham bajo la misma geometría.
El cambio de B1 a B16 no es idéntico bit a bit en todos los logits, aunque no
cambió ninguna decisión greedy en el diagnóstico fijado. Esas identidades y
límites permanecen registrados; no se generaliza la calificación CPU a GPU o XLA.
La evaluación de competencia citada procede de un cierre de recuperación;
los fallos operacionales anteriores permanecen en la procedencia histórica.

\subsection{Qué no se ha demostrado}
\begin{itemize}
\item Las dos baterías son repetición numérica de la misma población, no dos
semillas o mundos independientes. No justifican significancia estadística.
\item La ablación de un modelo entrenado mide acceso/predicción. No demuestra
que los hechos nunca se adquirieron o que una ruta los almacene de forma exclusiva.
\item La medida de NLL no equivale a clasificar una respuesta generada como
alucinación. La medición de canales no ejecutó ramas de entrenamiento H/T/C/R.
\item Una contribución local activa o un gradiente agregado activo no prueban
aprendizaje factual por unidad. Tampoco el efecto medio pequeño de U1 prueba muerte.
\end{itemize}

\callout{Hallazgo útil, hipótesis causal abierta}{
La evidencia vincula una intervención de acceso con cambios medibles de probabilidad
factual y muestra el límite del control original. La relación entre muerte de un
canal durante aprendizaje y afirmaciones falsas sigue sin probarse.}

\pageend
\section{V6: bloqueo persistente y distribuido}
\label{sec:v6}

El nuevo protocolo preserva \status\ y cambia explícitamente el estimando.
Compara políticas de bloqueo persistente/concentrado con interrupciones
distribuidas, manteniendo balance marginal de sitios e inputs en cada update
del panel. No intenta fabricar un match de U0 ni supone igual daño general
durante las trayectorias.

\figurepanel{prospective-design}{Calendarios derivados del plan V6 fijado.
\textbf{Es un diseño prospectivo: las ramas no se han entrenado.} Los colores
identifican el sitio cuyo veto se programa; no representan pérdidas o resultados
medidos. Los rescates retiran el bloqueo desde el índice de update 2000.}{fig:v6}

El panel tiene 17 ramas: H sano con sham, cuatro P con una unidad fija bloqueada,
cuatro D con bloqueos distribuidos y ocho rescates RP/RD. Todas parten del mismo
A0 dentro del mundo/seed y terminan a A4000 con los mismos ejemplos y horizonte.

Para $b=\lfloor t/4\rfloor$, $k=t\bmod4$ y una permutación $\pi_b$ fijada por
hash a partir de un seed independiente, sin consumir el RNG del modelo o los datos, D$_s$ veta
$\pi_b[(k+s)\bmod4]$. Cada conjunto de cuatro P o D asigna cada sitio exactamente
una vez en cada update. Cada D visita todos los sitios una vez por bloque de
cuatro; no queda ligada permanentemente al mismo slice auxiliar del currículo.

Cada P concentra 4.000 bloqueos programados en un sitio; cada D distribuye 1.000
por sitio. El balance no iguala COMMIT efectivos, energía borrada o gradientes
cuando divergen los parámetros. RP/RD coinciden con sus P/D hasta A2000 y usan
sham desde el update de índice 2000, que produce A2001. La inferencia primaria
es sin veto. Las ejecuciones son completas e independientes; no se cambia de
schedule mediante resume estricto.

\pageend
\section{Qué estimará V6 y qué hace falta para ejecutarlo}
\label{sec:future}

Para cada resultado final $Y$, las barras representan la media de las cuatro
ramas de cada familia:
\begin{align*}
\tau&=\overline{Y}_P-\overline{Y}_D, &
\rho_P&=\overline{Y}_{RP}-\overline{Y}_P,\\
\rho_D&=\overline{Y}_{RD}-\overline{Y}_D, &
\kappa&=\rho_P-\rho_D.
\end{align*}
$\tau$ es el contraste total entre políticas. $\kappa$ mide cuánto cambia su
diferencia al retirar ambos tipos de bloqueo. H contextualiza competencia y
deterioro; no es necesario asumir que los efectos de sitios distintos sean aditivos.

\subsection{Resultados, datos y límites}
El protocolo fija una realización nueva del mundo, exposición, modelo y calendario.
La unidad de replicación es el mundo/seed, no las 17 ramas, hechos o posiciones.
Se mantienen todos los fallos y no se reemplazan seeds por resultados desfavorables.

La evaluación primaria usa todos los hechos comunes de validation a A4000,
incluidos los que H falle; se reportan también raros y no expuestos. Las métricas
factuales promedian por hecho. La NLL original usa un inventario global distinto:
todas las ocurrencias TRAIN originales de entidades validation, ponderadas por
tokens válidos/EOS dentro de cada mapa. No se mezclan sus denominadores.

Los 17 entrenamientos se cierran y verifican antes de capturar validation.
Primero se sellan las generaciones y logits completos en la frontera pública,
sin el Value correcto; después se abren targets privados separados por split.
El lector histórico que carga un oráculo conjunto requiere sustitución calificada.
Validation del mundo anterior ya se evaluó en V2: no se presenta como globalmente
ciega por no haberse leído durante V5. Test permanece cerrado en el nuevo protocolo.

Se fija un margen descriptivo de 0,02 nats para interpretar conjuntamente NLL
original y deterioro factual por mapa. No es una prueba de equivalencia funcional
ni un gate que permita excluir, reponderar o repetir ramas después de observarlas.
Las curvas y diagnósticos de acceso están fijados de antemano. Un rescate finito
puede combinar recuperación de acceso y aprendizaje adicional; no identifica
automáticamente un mecanismo único ni prueba irreversibilidad cuando falla.

\subsection{Recursos y preparación técnica}
\begin{tabularx}{\linewidth}{@{}Y>{\raggedright\arraybackslash}p{48mm}@{}}
\toprule
\textbf{Concepto} & \textbf{Base histórica serial CPU}\\
\midrule
17 adquisiciones & 39 h 2 min (aprox.)\\
17 ramas con verificación, tres mapas y postrun & 56 h 18 min 50 s\\
Nuevo W640 e inicialización compartida & +14 min 30 s (aprox.)\\
Ejemplo de pre/post completos únicamente & 398,4375 GiB para 17 ramas\\
\bottomrule
\end{tabularx}

Estas cifras son extrapolaciones de terminales históricos, no un cap suficiente
del runner nuevo. Excluyen calificación, nuevas curvas, diagnósticos, otras
capturas, copias y contingencia. No se descuentan paralelización o prefijos
compartidos. GPU/AMD requiere calificación y medición de coste propias.

El calendario pasó 62 pruebas y el paquete seis comprobaciones adicionales. Eso
valida componentes puros de diseño. Todavía faltan la nueva familia de schedules
en kernel/guard/Native, el export y lector privado, inicialización y ruta numérica
calificadas, además del presupuesto completo. No hubo entrenamiento V6.

\pageend
\section{Reproducibilidad y procedencia}
\label{sec:provenance}

Este informe se distribuye con el código LaTeX, las cinco figuras vectoriales,
datos reducidos y scripts de reconstrucción. El objetivo es que un clon limpio
pueda reconstruir las gráficas y recalcular la selección, sin cargar un modelo,
consultar Native o descargar las capturas raw completas.

\begin{tabularx}{\linewidth}{@{}p{43mm}Y@{}}
\toprule
\textbf{Artefacto incluido} & \textbf{Qué permite comprobar}\\
\midrule
\texttt{data/report-data.json} & Datos curados de competencia, unidades, posiciones, hechos, límites y calendario ilustrado.\\
\texttt{data/provenance.json} & Fuentes originales, campos de extracción y hashes; diferencia entre originales y exportaciones reducidas.\\
Tablas y selector en \texttt{data/} & Recalcular agregados y decisión de ambas baterías bajo la regla original, desde sus tablas completas.\\
\texttt{design/} & Protocolo V6, calendario completo y comprobación pura del balance.\\
\texttt{make\_figures.py} & Reconstruir figuras desde los datos portables; sin selección de ejemplos posterior.\\
\texttt{report.tex} y \texttt{Makefile} & Compilar el PDF técnico; no iniciar experimentos.\\
\bottomrule
\end{tabularx}

La reproducción de tablas y figuras no sustituye la revalidación completa desde
los archivos raw, checkpoints y ledgers originales, que no se incluyen en este
paquete. El contador del readback histórico registró aproximadamente 19,66 GB
leídos; no es el tamaño conjunto de esos tres tipos de evidencia. El exportador
autentica los JSON y documentos fuente fijados cuando están disponibles, sin
revalidar los blobs raw ni los checkpoints. El paquete tampoco reconstruye la trayectoria
de entrenamiento histórica ni califica una ejecución futura por contener sus hashes.

\subsection{Identificadores principales}
Los paths repo-relative y la lista completa de fuentes están en el manifiesto
de procedencia. Los siguientes SHA256 identifican cierres y reglas citados:

{\footnotesize
\begin{tabularx}{\linewidth}{@{}p{32mm}Y@{}}
\toprule
\textbf{Evidencia} & \textbf{SHA256 del artefacto original}\\
\midrule
Competencia V5 & \texttt{bdb63840ae28176e5cbf22aded90ad78585804415e0d5579ca9bb63da1bee9db}\\
Cierre worker & \texttt{f223ec5ec1101a67528a4cd54d38ce9e4d7dd405ad0c26586f63c360d42810a2}\\
Cierre supervisor & \texttt{18b41a1980e6eef486a1747789cbb04b06154f403da7d52de9780a968acc6002}\\
Selección V5 & \texttt{ef1b3cbb783364b662974c0a605b42a9758433b2b77c864a644d08f1d6bfa78f}\\
Regla V1 & \texttt{a9f1aef9f00e5bb98aa7fe9e294771f279c0693a9815ae40a1ea34167d469f79}\\
Protocolo V6 & \texttt{6bae2857c9bf6910aa728be2cd65da5f510400190f2e45a1d909194bd48ea195}\\
Plan V6, archivo & \texttt{feb9d0cbb236831a864a4f8981772262953f198fa53570ecde01b3704c95f766}\\
\bottomrule
\end{tabularx}}

La fuente causal congelada de la medición corresponde al commit:
\begin{center}\texttt{033450e50d6a13b355b010cabaacf686a741d25f}\end{center}
Se conservan las identidades
separadas de esa fuente y de la referencia histórica. Ningún hash de este informe
confiere admisión científica o autorización de entrenamiento.

\subsection{Lectura final}
Los resultados justifican medir la contribución factual además de observar
actividad y pérdida agregada. También documentan un límite real del diseño
inicial: no existe control admisible entre los tres candidatos bajo la regla
fijada. V6 propone una comparación distinta y explícita para estudiar la
persistencia del bloqueo. La hipótesis causal permanece abierta hasta ejecutar
y evaluar ese estudio dentro de su alcance declarado.

\end{document}
